% EDITA AQUÍ: remitente y asunto; edita destinatario, cuerpo y firma en sections/.
\documentclass[11pt,a4paper]{article}
\usepackage[T1]{fontenc}
\usepackage[spanish,es-nodecimaldot,es-noquoting]{babel}
\usepackage{tgpagella}
\usepackage[margin=26mm,top=23mm,bottom=25mm]{geometry}
\usepackage{xcolor,fancyhdr}
\definecolor{accent}{HTML}{245C61}
\definecolor{muted}{HTML}{59676B}
\usepackage[colorlinks=true,allcolors=accent]{hyperref}
\newcommand{\sendername}{Tu nombre}
\newcommand{\senderrole}{Tu cargo / Organización}
\newcommand{\senderemail}{you@example.com}
\newcommand{\lettersubject}{Asunto de la carta}
\hypersetup{pdftitle={\lettersubject},pdfauthor={\sendername}}
\setlength{\parindent}{0pt}
\setlength{\parskip}{10pt}
\linespread{1.08}
\pagestyle{fancy}\fancyhf{}
\fancyfoot[L]{\small\color{muted}\sendername\quad /\quad\href{mailto:\senderemail}{\senderemail}}
\fancyfoot[R]{\small\color{muted}\thepage}
\renewcommand{\headrulewidth}{0pt}
\begin{document}
{\LARGE\bfseries\sendername}\par
\vspace{-3pt}{\small\sffamily\color{accent}\senderrole}\par
\vspace{2pt}
{\small\color{muted}Dirección, Ciudad, Código postal\\
\href{mailto:\senderemail}{\senderemail}\quad /\quad +00 123 456 789}
\vspace{8pt}{\color{accent}\hrule height 1.3pt}
\vspace{18pt}

\begin{minipage}[t]{0.62\linewidth}
\textbf{Nombre del destinatario}\\
Cargo / Departamento\\
Organización\\
Dirección\\
Ciudad, Código postal
\end{minipage}\hfill
\begin{minipage}[t]{0.32\linewidth}\raggedleft
{\small\color{muted}\today}
\end{minipage}
\vspace{22pt}


{\large\bfseries\color{accent}\lettersubject}

Estimado/a destinatario/a:

Expón el propósito de la carta en una frase clara. Si respondes a una oferta,
propuesta o conversación previa, identifícala aquí y explica por qué te diriges
a esta persona.

Desarrolla la idea principal con un ejemplo concreto. En una candidatura,
relaciona tu experiencia con las necesidades de la organización. En una solicitud
formal, aporta el contexto, la información de apoyo y el resultado que buscas.
Céntrate en lo que necesita el destinatario para tomar una decisión.

Añade otro párrafo de apoyo si es necesario. Usa pruebas concretas y dedica
cada párrafo a una idea. Menciona los documentos adjuntos y explica cómo
respaldan tu solicitud.

Cierra con un siguiente paso práctico: propone una reunión, solicita respuesta
para una fecha razonable o indica cómo contactar contigo. Agradece el tiempo
y la atención dedicados.

\vspace{5pt}
Atentamente,

\vspace{18pt}
\textbf{\sendername}\\
{\small\color{muted}\senderrole}

\vfill
{\small\color{muted}\textbf{Adjuntos:} CV / Documento de apoyo (elimina si no procede).}

\end{document}
